\subsection{Limitations}
\label{sec:conclusion-limitations}

While COSMO satisfies the AI-native criteria outlined above,
several limitations should be acknowledged. Each item below concerns a
feature that is implemented but falls short of its design; capabilities that
were never in scope are extensions and are listed under Future Work
(Section~\ref{sec:future-work}).
\label{limit}
\begin{itemize}
  \item \textbf{Retrieval is dense-only, and its cut-off is set rather than
    tuned.} Grounding uses a single nearest-neighbour search over embeddings:
    there is no keyword or BM25 leg to catch an exact term the embedding blurs
    (a product name, a plan name, a figure), and no reranking stage over the
    retrieved set. A chunk is admitted if its cosine distance from the query is
    0.5 or less, which keeps chunks at a similarity of 0.5 or better. That
    number was chosen to be defensibly strict rather than measured: choosing it
    properly needs a labelled retrieval set --- queries paired with the chunks
    that ought to answer them --- which this project does not have. The
    consequence is that retrieval quality is asserted from the pipeline's
    construction rather than measured, and it is the clearest gap between how
    the reply-drafting capability is evaluated (Section~\ref{sec:llm-judge})
    and how its inputs are.

    A related consequence is worth stating because it is easy to miss: the
    cut-off was previously 0.8, described in the code as ``fairly relevant''.
    A distance of 0.8 is a similarity of 0.2, so in practice the filter
    admitted almost anything and the five nearest chunks were injected into
    every prompt regardless of whether they bore on the question. A grounding
    step that always fires is not grounding; it is noise with a citation. The
    unit confusion behind that --- treating a distance as a similarity --- is
    the kind of defect a retrieval evaluation would have caught immediately
    and inspection did not.

  \item \textbf{Chunks are fixed-size and boundary-blind.} Documents are split
    every 1,000 characters with a 150-character overlap, without regard to
    headings, tables, or sentence ends. A pricing table split across two chunks
    is retrieved as two fragments, and the model may answer from the half it
    was given. Structure-aware splitting, and retrieving a chunk's neighbours
    alongside it, are the standard remedies and are not implemented.

  \item \textbf{Follow-ups fire on the no-reply window, not on the
    day-based cadence.} Section~4.4 describes Follow-up~1 on days~4--5
    and Follow-up~2 on days~9--12, but both steps currently fire on the
    same no-reply window. The day-based fields are stored, validated and
    editable by an administrator, and the settings page says on the fields
    themselves that the scheduler does not yet read them.

  \item \textbf{Daily action ranking is a single-shot LLM
    re-rank over a heuristic shortlist.} The model reorders the
    day's candidate actions in one batched call, but it does not
    learn from which actions the representative actually completed;
    incorporating that outcome signal is left as future work.

  \item \textbf{The HubSpot import is not functional.} Contact import
    works through Apollo.io, CSV upload, a LinkedIn Chrome Extension and
    inbound lead forms. A HubSpot import endpoint also exists, but it records
    the request and returns success without enqueueing any work, so no
    contacts are actually pulled.

  \item \textbf{The AI reply intent classifier uses an
    English-language prompt}; classification accuracy on
    non-English inbound replies has not been formally evaluated.

  \item \textbf{Email sending is tied to the representative's Gmail
    account.} Sending goes through Gmail OAuth, so deliverability depends on
    that account's sending reputation and on Gmail's rate limits.

\end{itemize}

\subsection{Future Work}
\label{sec:future-work}

The limitations above are engineering work on paths that already exist ---
the HubSpot import can be wired to the existing contact-import worker, and the
follow-up scheduler can read the cadence fields it already stores. This section is about extending COSMO beyond its current scope, in
four groups.

\textbf{New capabilities on the existing platform.} Password authentication
alongside Google OAuth, for organisations that restrict third-party sign-in;
dedicated outbound mailboxes (custom SMTP or a provider such as SendGrid) so
that sending no longer depends on one Gmail account; and a proactive Ask COSMO,
which opens each morning with the overnight replies and the follow-ups due. The
last needs no new data: the daily action list and the notification channel
already exist, and it would be bounded by the same relevance threshold,
frequency ceiling and off switch as automatic replies.

\textbf{Measuring whether the AI helps.} The system tracks reply rate and
meetings booked, but cannot yet show that its generated content outperforms a
human-written email. That requires open and click tracking, an A/B framework over comparable contact cohorts, and enough
volume for a difference to be measurable; the feedback tables already record
whether each draft was used, edited or discarded. Recording which daily actions
were completed, against the rank they were given, would likewise let the ranking
be evaluated and tuned rather than asserted.

\textbf{Extending the evaluation.} The evaluation covers English mail only;
Vietnamese replies need their own labelled seed set, since intent often lives in
politeness register that translation flattens. More than two human raters would
let agreement be reported with a confidence interval. For the classifier, the
concrete target is the wrong-recipient scenario on which three of its four
misses occur (Section~\ref{sec:llm-judge}): a prompt revision for that case, then
distillation from a stronger model. Retrieval also needs a labelled query set so
that its relevance cut-off can be tuned rather than chosen.

\textbf{Directions for the Graduation Project.} LLM-based ICP scoring over the
contact's full profile and interaction history, replacing the weighted formula;
integration with Salesforce and LinkedIn Sales Navigator and a two-way
HubSpot sync, so that COSMO acts as the outreach layer over a pipeline a team
already maintains; and multi-channel orchestration
beyond email, including LinkedIn and, for the Vietnamese market, Zalo.


\clearpage
